\scanpage{108}
\mspage{53}
% archive sheet number: 107
pas, en effet, qu'il y en ait -- la difficulté, par cette voie, paraît irrémédiable. Par contre, si on regarde $\Phi',\Psi'$ comme foncteurs à valeurs dans $\underline{\mathrm{Fib}}_0\,S,\underline{\mathrm{Cofib}}_0\,S$ respectivement (donc dans des catégories où on n'exige pas que les flèches soient cartésiennes resp. cocartésiennes), on trouve encore des $W$-morphismes

% slanted
\margin{Mais c'est idiot de prendre, de \unsure{pareils|nouveaux} foncteurs si compliqués, \unsure{parce|pour} qu'on veut tout ramener à $\mathrm{Cat}_{/S}$. Il \unsure{serait} plus simple de comparer $\underline{\mathrm{Fibhot}}\,S$ directement à $\underline{\mathrm{Fibsc}}_0\,S$ \deleted{\unsure{etc.}} \unsure{à} $\underline{\mathrm{Cofibsc}}_0\,S$ \unsure{via|avec} $\Phi_0$ et $\Psi_0$ (qui sont en fait définis sur $\underline{\mathrm{Fibhot}}\,S$, pas seulement sur $\underline{\mathrm{Parf}}$) ($\Psi_0$ et $\underline{\mathrm{Cofibsc}}_{\unsure{0}/S}\,S$), et on trouve comme ça (à des sens près) $\mathrm{id}\deleted{\unsure{*}}\xrightarrow{W}\pi_0\Phi_0$ et $\Phi_0\pi_0\xrightarrow{W}\mathrm{id}$ où $\pi:\underline{\mathrm{Fibsc}}_0\,S\longrightarrow\underline{\mathrm{Fibhot}}\,S$.}

\[
 \mathrm{id}\longrightarrow\Phi'p,\qquad \mathrm{id}\longrightarrow\Psi'q,
\]

ce qui liquide au moins déjà le cas de toutes les catégories du diagramme \deleted{\unsure{que}} p.~8 qui sont des sous-catégories pleines de $\mathrm{Cat}_{/S}$. Il ne reste \deleted{que} à traiter que le cas des quatre catégories de structures ``scindées''

\[
 \underline{\mathrm{Fibsc}}\,S,\ \underline{\mathrm{Fibsc}}_0\,S \quad\text{et}\quad \underline{\mathrm{Cofibsc}}\,S,\ \underline{\mathrm{Cofib}}_0\,S.
\]

Bien sûr il suffit de le faire d'un côté, disons du côté des fibrés, par dualité ($S\mathbin{\unsure{\mapstochar\relbar|\longmapsto}}S^o$) cela établira le résultat symétrique. On va donc regarder d'un peu plus près le carré de foncteurs

\diagram{\begin{tikzcd}[ampersand replacement=\&,column sep=large,row sep=large]
 \underline{\mathrm{Fibsc}}\,S\arrow[d,"p"] \& \underline{\mathrm{Fibsc}}_0\,S\arrow[l,hookrightarrow,"\alpha'"']\arrow[d,"p_0","{\unsure{\prime\prime}}"']\\
 \underline{\mathrm{Fib}}\,S\arrow[u,bend left=36,dashed,"\varphi"] \& \mathrm{Fib}_{0\deleted{\unsure{\mathrm{sc}}}}\,S\arrow[l,hookrightarrow,"\alpha"']\arrow[u,bend right=36,dashed,"\varphi_0"']
\end{tikzcd}}

% slanted
\margin{NB Il s'agit de prouver seulement que $p,p_0$ induisent des équiv. pour les localisés (pour la $W$-équiv. grossière). On sait déjà qu'il en est ainsi pour $\alpha$, et ça en résultera pour $\alpha'$ !}

% slanted
\margin{inutile de regarder ce carré, cf annotation \unsure{marginale} ci-dessus (haut de la page)}

le foncteur

\[
 \Phi:\mathrm{Cat}_{/S}\longrightarrow\underline{\mathrm{Fibsc}}\,S
\]

induit des foncteurs en sens inverse de $p,p_0$, sur les sous-catégories pleines $\underline{\mathrm{Fib}}\,S,\mathrm{Fib}_0\,S$ de $\mathrm{Cat}_{/S}$, foncteurs marqués en pointillés comme $\varphi,\varphi_0$.

